% Options for packages loaded elsewhere
\PassOptionsToPackage{unicode}{hyperref}
\PassOptionsToPackage{hyphens}{url}
\documentclass[
  a4paper,
]{article}
\usepackage{xcolor}
\usepackage[margin=25mm]{geometry}
\usepackage{amsmath,amssymb}
\setcounter{secnumdepth}{-\maxdimen} % remove section numbering
\usepackage{iftex}
\ifPDFTeX
  \usepackage[T1]{fontenc}
  \usepackage[utf8]{inputenc}
  \usepackage{textcomp} % provide euro and other symbols
\else % if luatex or xetex
  \usepackage{unicode-math} % this also loads fontspec
  \defaultfontfeatures{Scale=MatchLowercase}
  \defaultfontfeatures[\rmfamily]{Ligatures=TeX,Scale=1}
\fi
\usepackage{lmodern}
\ifPDFTeX\else
  % xetex/luatex font selection
\fi
% Use upquote if available, for straight quotes in verbatim environments
\IfFileExists{upquote.sty}{\usepackage{upquote}}{}
\IfFileExists{microtype.sty}{% use microtype if available
  \usepackage[]{microtype}
  \UseMicrotypeSet[protrusion]{basicmath} % disable protrusion for tt fonts
}{}
\makeatletter
\@ifundefined{KOMAClassName}{% if non-KOMA class
  \IfFileExists{parskip.sty}{%
    \usepackage{parskip}
  }{% else
    \setlength{\parindent}{0pt}
    \setlength{\parskip}{6pt plus 2pt minus 1pt}}
}{% if KOMA class
  \KOMAoptions{parskip=half}}
\makeatother
\usepackage{longtable,booktabs,array}
\newcounter{none} % for unnumbered tables
\usepackage{calc} % for calculating minipage widths
% Correct order of tables after \paragraph or \subparagraph
\usepackage{etoolbox}
\makeatletter
\patchcmd\longtable{\par}{\if@noskipsec\mbox{}\fi\par}{}{}
\makeatother
% Allow footnotes in longtable head/foot
\IfFileExists{footnotehyper.sty}{\usepackage{footnotehyper}}{\usepackage{footnote}}
\makesavenoteenv{longtable}
\ifLuaTeX
  \usepackage{luacolor}
  \usepackage[soul]{lua-ul}
\else
  \usepackage{soul}
\fi
\setlength{\emergencystretch}{3em} % prevent overfull lines
\providecommand{\tightlist}{%
  \setlength{\itemsep}{0pt}\setlength{\parskip}{0pt}}
\usepackage{newunicodechar}
\newunicodechar{↔}{\ensuremath{\leftrightarrow}}
\newunicodechar{→}{\ensuremath{\rightarrow}}
\newunicodechar{←}{\ensuremath{\leftarrow}}
\newunicodechar{≥}{\ensuremath{\ge}}
\newunicodechar{≤}{\ensuremath{\le}}
\newunicodechar{×}{\ensuremath{\times}}
\newunicodechar{≪}{\ensuremath{\ll}}
\newunicodechar{≫}{\ensuremath{\gg}}
\newunicodechar{∈}{\ensuremath{\in}}
\newunicodechar{∞}{\ensuremath{\infty}}
\newunicodechar{±}{\ensuremath{\pm}}
\newunicodechar{−}{\ensuremath{-}}
\newunicodechar{≈}{\ensuremath{\approx}}
\newunicodechar{≠}{\ensuremath{\ne}}
\newunicodechar{∝}{\ensuremath{\propto}}
\newunicodechar{√}{\ensuremath{\sqrt{\phantom{x}}}}
\newunicodechar{⊕}{\ensuremath{\oplus}}
\newunicodechar{α}{\ensuremath{\alpha}}
\newunicodechar{β}{\ensuremath{\beta}}
\newunicodechar{θ}{\ensuremath{\theta}}
\newunicodechar{ρ}{\ensuremath{\rho}}
\newunicodechar{π}{\ensuremath{\pi}}
\newunicodechar{φ}{\ensuremath{\varphi}}
\newunicodechar{λ}{\ensuremath{\lambda}}
\newunicodechar{σ}{\ensuremath{\sigma}}
\newunicodechar{Δ}{\ensuremath{\Delta}}
\newunicodechar{Σ}{\ensuremath{\Sigma}}
\newunicodechar{ν}{\ensuremath{\nu}}
\newunicodechar{κ}{\ensuremath{\kappa}}
\newunicodechar{μ}{\ensuremath{\mu}}
\newunicodechar{τ}{\ensuremath{\tau}}
\newunicodechar{ω}{\ensuremath{\omega}}
\newunicodechar{ε}{\ensuremath{\varepsilon}}
\newunicodechar{Ψ}{\ensuremath{\Psi}}
\newunicodechar{⁻}{\ensuremath{{}^{-}}}
\newunicodechar{¹}{\ensuremath{{}^{1}}}
\newunicodechar{²}{\ensuremath{{}^{2}}}
\newunicodechar{³}{\ensuremath{{}^{3}}}
\newunicodechar{⋆}{\ensuremath{\star}}
\newunicodechar{∗}{\ensuremath{*}}
\newunicodechar{★}{\ensuremath{\bigstar}}
\newunicodechar{∎}{\ensuremath{\blacksquare}}
\newunicodechar{ℝ}{\ensuremath{\mathbb{R}}}
\newunicodechar{ℤ}{\ensuremath{\mathbb{Z}}}
\newunicodechar{₀}{\ensuremath{{}_{0}}}
\newunicodechar{₁}{\ensuremath{{}_{1}}}
\newunicodechar{₂}{\ensuremath{{}_{2}}}
\newunicodechar{₃}{\ensuremath{{}_{3}}}
\newunicodechar{₄}{\ensuremath{{}_{4}}}
\newunicodechar{₅}{\ensuremath{{}_{5}}}
\newunicodechar{₆}{\ensuremath{{}_{6}}}
\newunicodechar{₇}{\ensuremath{{}_{7}}}
\newunicodechar{₈}{\ensuremath{{}_{8}}}
\newunicodechar{₉}{\ensuremath{{}_{9}}}
\newunicodechar{✓}{\ensuremath{\checkmark}}
\newunicodechar{—}{---}
\newunicodechar{①}{(1)}
\newunicodechar{②}{(2)}
\newunicodechar{③}{(3)}
\usepackage{bookmark}
\IfFileExists{xurl.sty}{\usepackage{xurl}}{} % add URL line breaks if available
\urlstyle{same}
\hypersetup{
  hidelinks,
  pdfcreator={LaTeX via pandoc}}

\author{}
\date{}

\begin{document}

\section{Two-Grammar Decomposition of Interaction in an Anonymous
Two-Channel Closed Wave
System}\label{two-grammar-decomposition-of-interaction-in-an-anonymous-two-channel-closed-wave-system}

\subsection{A Classification Theorem for Conserved Readouts,
Term-by-Term Machine-Precision Verification of the Two-Term Flow
Identity, and the Necessary Bifurcation of Gravity-Type and Charge-Type
Composition Laws from Wave
Shape}\label{a-classification-theorem-for-conserved-readouts-term-by-term-machine-precision-verification-of-the-two-term-flow-identity-and-the-necessary-bifurcation-of-gravity-type-and-charge-type-composition-laws-from-wave-shape}

\textbf{Author:} Noriaki Kihara\\
\textbf{ORCID:} 0009-0004-6753-4020\\
\textbf{Date:} August 2, 2026\\
\textbf{Version DOI:} \texttt{10.5281/zenodo.21763996}\\
\textbf{Concept DOI:} \texttt{10.5281/zenodo.21763995}\\
\textbf{Position:} ``Wave Information Readout'' series, two-channel
exchange-scattering system, additional paper v1 (paper-9 series;
\textbf{Part I of a trilogy}. Part II = Structure of Counting Readouts;
Part III = Lock Dynamics. For the division, see ``Position of This Paper
and the Three-Part Structure'' below.)

\begin{center}\rule{0.5\linewidth}{0.5pt}\end{center}

\subsection{Abstract}\label{abstract}

\textbf{(Discovery of the mechanism)} The two-body flow in an anonymous
closed wave system decomposes identically into

\[
dN_B
\;=\;
\underbrace{\sin^2\theta\,(N_A - N_B)}_{\text{magnitude term}}
\;+\;
\underbrace{\sin 2\theta\,\mathrm{Re}\langle a|b\rangle}_{\text{overlap term}}
\]

\textbf{two terms, and no third term exists}. These two terms read,
respectively and \textbf{exclusively}, the two kinds of memory carried
by the shape of a wave --- the norm (magnitude), and the relative phase
on shared modes (shape). Consequently, in a wave system possessing no
labels, no background, and no external fields, interaction necessarily
differentiates into exactly \textbf{two kinds}:

{\def\LTcaptype{none} % do not increment counter
\begin{longtable}[]{@{}
  >{\raggedright\arraybackslash}p{(\linewidth - 4\tabcolsep) * \real{0.3333}}
  >{\raggedright\arraybackslash}p{(\linewidth - 4\tabcolsep) * \real{0.3333}}
  >{\raggedright\arraybackslash}p{(\linewidth - 4\tabcolsep) * \real{0.3333}}@{}}
\toprule\noalign{}
\begin{minipage}[b]{\linewidth}\raggedright
\end{minipage} & \begin{minipage}[b]{\linewidth}\raggedright
Magnitude term
\end{minipage} & \begin{minipage}[b]{\linewidth}\raggedright
Overlap term
\end{minipage} \\
\midrule\noalign{}
\endhead
\bottomrule\noalign{}
\endlastfoot
Composition law & \textbf{Additive} (mean, error
\(\le7.2\times10^{-16}\)) & \textbf{Product} \(\sqrt{f_Af_B}\) (error
\(\le4.0\times10^{-15}\)) \\
Sign & None (only the norm difference sets the direction) & Present
(reverses with relative phase, measured) \\
Screening & Impossible (always on whenever a norm is present) & Possible
(two mechanisms: zero component / non-shared channels) \\
Quantization & None (continuous) & Present (rational locking,
measured) \\
Range of the coupling & Ratio \((R/R_0)^2\): unbounded, continuous
(\textbf{weakness by orders of magnitude is possible}) & Angle:
unitarity bound \(\alpha\le1/4\pi\) × lock discretization
(\textbf{confined to the \(O(10^{-1})\) band}) \\
\end{longtable}
}

This property table \textbf{corresponds one-to-one with the contrast
table of gravity (additive mass, no sign, unscreenable, unquantized) and
electromagnetism (product of charges, ±, screenable, quantized)}. This
correspondence is not an arbitrary assignment based on resemblance of
properties: if the overlap term is read as gravity-type, gravitational
action would reverse and vanish under relative-phase reversal and
channel non-sharing, contradicting universality; if the magnitude term
is read as charge-type, none of the amplitude product, the two signs, or
neutralization can be constructed --- \textbf{the assignment is a unique
correspondence that excludes the reversed assignment} (Section 6.5,
Corollary). That is, this paper discovers the mechanism by which the
divergence in the properties of the two great long-range interactions
arises necessarily, from within a single identity, as the exclusive
readout of the two memories of wave shape, and verifies every row at
machine precision. Anonymity is never broken --- distinguishing the
forces requires neither particle species nor external input of coupling
constants.

\textbf{(Completeness on the Coulomb side)} The overlap term possesses
all five elements of the Coulomb structure --- inverse square (already
derived from harmonic closure {[}3{]}), product of couplings, sign,
quantization, and value (the exact finite-order root
\(\sin^2(23\pi/124)\), matching the dimensionless elementary charge
\(\sqrt{4\pi\alpha}\) to 0.16 ppm {[}2{]}). \textbf{We call this
quantity the model's elementary charge.}

\textbf{(Hierarchy)} A long-range interaction that is weaker by orders
of magnitude can exist only in the grammar of the magnitude term.
Moreover, the coupling of the magnitude term is of scale-ratio type
\((R/R_0)^2\), and in any universe containing localized closures
\(R\ll R_0\) --- \textbf{the enormous hierarchy between the two great
long-range forces is not a coincidence requiring explanation but a
structural necessity that follows from the very definition of
localization: particles are far smaller than the background}. What
remains underived is not the existence of the hierarchy but only its
exponent (the scale dictionary).

\textbf{(The one remaining item)} Toward identification with the
physical elementary charge, the only work remaining \textbf{inside the
model} is the derivation of the mechanism that selects order 124. The
static decomposition does not select it (refuted); the fixed-point
channel has been eliminated; the lock weights collapse with the
denominator (measured ratio \(>461\)); and observational selection
passes the selection down to the divisor structure of the observation
clock --- the selection problem has been transformed into a single
falsifiable question, ``why is the observation clock commensurable with
248?'' (sister paper C).

\textbf{(Convergence)} The force law product × cos(relative phase) ×
inverse square was realized in fluids by Bjerknes (1875, pulsating
spheres); mode overlap as the existence condition of coupling is
realized by the visibility law of optics; quantization plateaus via
locking by the Shapiro steps of superconductivity; and the
correspondence between force ratios and scale ratios by Dirac's large
numbers (1937) --- each independently realized and observed. That
mutually unrelated systems converge on the same grammar is evidence that
this grammar is not a detail of the medium but a general structure of
closed wave systems. The contribution specific to this paper is to show
that, in an endogenous-coupling system in which the coupling is not
supplied externally, all five elements and every row of the two grammars
are derived and verified simultaneously as terms of a single flow
equation, and survive an audit over five variants of the readout
convention.

\textbf{Keywords:} two-channel exchange scattering, conserved readouts,
composition laws, interference cross term, Bjerknes force, Arnold
tongue, charge quantization, hierarchy problem, reproducible computation

\begin{center}\rule{0.5\linewidth}{0.5pt}\end{center}

\subsection{0. Conclusion}\label{conclusion}

\[
\boxed{
\begin{aligned}
&\text{The two-body flow of an anonymous closed wave system decomposes identically into a "magnitude term" and an "overlap term"; there is no third term.}\\
&\text{Magnitude term: additive composition (error }7.2\times10^{-16}\text{), no sign, unscreenable; its readout is exactly conserved.}\\
&\text{Overlap term: product composition }\sqrt{f_Af_B}\text{ (error }4.0\times10^{-15}\text{), sign = relative phase, neutralized by two mechanisms, quantized by rational locking.}\\
&\text{This contrast corresponds one-to-one with the property table of gravity/electromagnetism, and only the magnitude term can become weaker by orders of magnitude.}
\end{aligned}
}
\]

\begin{center}\rule{0.5\linewidth}{0.5pt}\end{center}

\subsection{Position of This Paper and the Three-Part
Structure}\label{position-of-this-paper-and-the-three-part-structure}

This work reports results obtained from a single experimental
environment, divided into three parts according to the independence of
the claims.

{\def\LTcaptype{none} % do not increment counter
\begin{longtable}[]{@{}
  >{\raggedright\arraybackslash}p{(\linewidth - 6\tabcolsep) * \real{0.2500}}
  >{\raggedright\arraybackslash}p{(\linewidth - 6\tabcolsep) * \real{0.2500}}
  >{\raggedright\arraybackslash}p{(\linewidth - 6\tabcolsep) * \real{0.2500}}
  >{\raggedright\arraybackslash}p{(\linewidth - 6\tabcolsep) * \real{0.2500}}@{}}
\toprule\noalign{}
\begin{minipage}[b]{\linewidth}\raggedright
\end{minipage} & \begin{minipage}[b]{\linewidth}\raggedright
Subject
\end{minipage} & \begin{minipage}[b]{\linewidth}\raggedright
Main results
\end{minipage} & \begin{minipage}[b]{\linewidth}\raggedright
Status
\end{minipage} \\
\midrule\noalign{}
\endhead
\bottomrule\noalign{}
\endlastfoot
\textbf{Part I (this paper)} & Two-grammar decomposition of interaction
& Two-term flow identity, classification theorem for conserved readouts,
composition laws (additive/product), sign, neutralization, direction of
the hierarchy & This paper \\
\textbf{Part II (sister paper B)} & Structure of counting readouts &
Exact rationality of equal-weight points, meta counting rule (bin table
× mask), full report of the convention audit, complete verification of
the Niven intersections (orders 3/4/6 = crystallographic menu), readout
device for the integer pair \((m,n)\) & Published (DOI:
10.5281/zenodo.21763998) \\
\textbf{Part III (sister paper C)} & Lock dynamics and the limits of
selection & Full report of the order-6 Arnold tongue, measured tongue
widths (ratio \(>461\)), fixed-point map (elimination of the selection
channel), observational selection (clock inheritance and
finite-observation limits), supercritical diagnostics & Published (DOI:
10.5281/zenodo.21764000) \\
\end{longtable}
}

The interdependence is as follows. Claims 1--3 of this paper stand on
their own; the results of sister papers B and C are cited only as a
concrete instance of the quantization row (summarized in Section 7) and
as the guarantee of convention robustness (summarized in Section 8).
Sister paper B shows in complete form that the claims of this paper do
not depend on the readout convention. Sister paper C carries the search
program for Claim 4 (the selection problem). All three parts share the
reproduction package of the same repository (Section 12).

\subsection{1. Research Question}\label{research-question}

In a system consisting only of waves without labels, how does
interaction differentiate into kinds? Without externally supplying
particle species, coupling constants, or background fields, can forces
of the distinct grammars --- gravity-type (universal, additive,
unsigned) and charge-type (product, both signs, quantized) --- arise
simultaneously from the same substrate? If they can, which information
in the wave shape carries the distinction, and can the exhaustiveness of
the distinction (the absence of a third type) be established?

In prior work, the operator structure and finite-order resonance of the
two-channel exchange-scattering system {[}2{]}, the inverse-square law
from harmonic closure {[}3{]}, and conserved-quantity readout via
interference {[}4{]} were each established separately. This paper
unifies them on a single flow equation and takes the mechanism of
differentiation itself as the main result.

\subsection{2. Definition of the System and Design
Boundary}\label{definition-of-the-system-and-design-boundary}

The system is the two-channel complex field
\(a(\chi,\eta), b(\chi,\eta)\) of the prior experimental environment
(the paper-9-series control-experiment / wave-packet-contraction
execution environment; repository directory romanized as
\texttt{taishou-jikken\_hasoku-shuushuku\_jikkou-kankyou\_v1/ab\_invariant\_theta\_toy\_v1}).
States use only the harmonic-packet construction via the existing
\texttt{explicit\_packet\_case} / \texttt{make\_case\_state}, and in the
standard configuration a carrier (\(q_A=+1, q_B=-1\)) is attached. The
readout \(\theta\) is defined from the power ratio of the even,
\(|k|\ge4\) sector of the joint \(\chi\) spectrum as

\[
\theta=\operatorname{atan2}\!\left(\sqrt{P_f},\,\sqrt{P_b}\right),
\qquad
P_f=\sum_{\text{mask}}\left(|A_k|^2+|B_k|^2\right)
\]

and one collision is the real orthogonal rotation

\[
a'=a\cos\theta-b\sin\theta,
\qquad
b'=a\sin\theta+b\cos\theta
\]

(the sign convention of the rotation was fixed to \(\sigma=+1\) by a
calibration experiment).

\textbf{Design boundary:} The scattering body and the \(\theta\)-readout
formula are unchanged across all experiments. Target values (fractions
\(f\), etc.) are used only for state construction and are never passed
to the forward processing. The variant readouts and probe rotations of
Section 6 are distinguished as explicitly declared working hypotheses /
instrument settings.

\textbf{Intrinsic transmission fraction:} The fraction \(f(s)\) of a
state \(s\) is defined as the in-mask power ratio of \(s\) alone. Via
the orthogonal two-harmonic mixture

\[
s(f)=\sqrt{1-f}\,s_{\mathrm{bos}}+\sqrt{f}\,s_{\mathrm{fer}}
\]

of the all-bosonic single harmonic (\(f=0\)) and the all-fermionic
single harmonic (\(f=1\)), \(f\) can be tuned continuously and exactly
(fraction error \(\le 10^{-31}\), cross overlap
\(\le 1.3\times10^{-16}\)).

\subsection{3. The Two-Term Flow
Identity}\label{the-two-term-flow-identity}

Direct computation for the real orthogonal rotation gives the
per-collision change of the B-channel norm as

\[
\boxed{
dN_B
=\underbrace{\sin^2\theta\,(N_A-N_B)}_{\text{magnitude term}}
+\underbrace{\sin 2\theta\,\mathrm{Re}\langle a|b\rangle}_{\text{overlap term}}
}
\]

This is an identity of the norm algebra of rotations, and \textbf{no
third term exists in this class of two-channel unitary flows}. The
magnitude term reads only norms (diagonal quantities); the overlap term
reads only the inner product \(\langle a|b\rangle\) (the amplitude
product and relative phase on shared modes). All experiments below
measure, term by term, the composition law, sign, conservation, and
quantization of these two terms.

Note that the vanishing of the magnitude term at \(N_A=N_B\) does not
mean the vanishing of the coupling; it means that the net transfer flow
balances between two symmetric channels. The coupling coefficient
\(\sin^2\theta\) and the observed net flow \(\sin^2\theta\,(N_A-N_B)\)
are distinct --- the same distinction as the fact that thermal coupling
between two bodies at equal temperature does not vanish even though the
net heat flow is zero.

For the mixed states \(a=s_A(f_A)\),
\(b=\nu\,[\sqrt{1-f_B}\,B_1+e^{i\varphi}\sqrt{f_B}\,B_7]\) (\(\nu\):
norm multiplier), the exact prediction is

\[
dN_B
=\sin^2\theta\,(1-\nu^2)
+\sin2\theta\,\nu\!\left[\sqrt{(1-f_A)(1-f_B)}\,c_0
+\sqrt{f_Af_B}\,c_6\cos\varphi\right]
\]

where \(c_0=\langle A_1|B_1\rangle\) and \(c_6=\langle A_5|B_7\rangle\)
are couplings measured from the state construction alone.

\subsection{4. Classification Theorem for Conserved
Readouts}\label{classification-theorem-for-conserved-readouts}

\textbf{Theorem (sufficient condition for conservation).} The real
orthogonal rotation exactly conserves \(|A_k|^2+|B_k|^2\) in each
frequency bin (the parallelogram identity). Hence any readout that is a
function only of the bin-wise rotation invariants is exactly conserved
under the evolution.

\textbf{Numerical classification.} Seven candidate \(\theta\) readouts
were classified by their self-drift under the dynamics that each readout
itself drives (300 collisions × 3 amplitudes).

{\def\LTcaptype{none} % do not increment counter
\begin{longtable}[]{@{}
  >{\raggedright\arraybackslash}p{(\linewidth - 4\tabcolsep) * \real{0.3000}}
  >{\raggedright\arraybackslash}p{(\linewidth - 4\tabcolsep) * \real{0.3000}}
  >{\raggedleft\arraybackslash}p{(\linewidth - 4\tabcolsep) * \real{0.4000}}@{}}
\toprule\noalign{}
\begin{minipage}[b]{\linewidth}\raggedright
Readout
\end{minipage} & \begin{minipage}[b]{\linewidth}\raggedright
Classification
\end{minipage} & \begin{minipage}[b]{\linewidth}\raggedleft
Max self-drift
\end{minipage} \\
\midrule\noalign{}
\endhead
\bottomrule\noalign{}
\endlastfoot
R1 diagonal sum \(\sum(\lvert A_k\rvert^2+\lvert B_k\rvert^2)\) &
\textbf{Conserved (unique)} & \(2.2\times10^{-16}\) \\
R2 \(\lvert A+B\rvert^2\) (X power) & Dynamical & \(1.06\) \\
R3 \(\lvert A-B\rvert^2\) & Dynamical & \(1.06\) \\
R4 cross spectrum \(\lvert\sum A_k B_k^*\rvert\) & Dynamical
(fixed-point convergence) & \(0.43\) \\
R5 \(\lvert A\rvert^2\) only / R6 \(\lvert B\rvert^2\) only & Dynamical
& \(1.35\) \\
R7 cross real part & Dynamical & \(0.38\) \\
\end{longtable}
}

Only the diagonal sum was conserved. This single structure
simultaneously generates the three no-gos described below --- the
freezing of \(\theta\), additive composition, and the absence of product
structure in the diagonal sector. Note that R4 was initially
misidentified as conserved on the basis of its tail standard deviation
of \(10^{-16}\), but the classification experiment revealed convergence
to a fixed-point attractor (the history of the correction is recorded in
the Reproducibility section).

\subsection{5. Experiment I: Composition Law of the Magnitude Term
(Additive, Gravity
Grammar)}\label{experiment-i-composition-law-of-the-magnitude-term-additive-gravity-grammar}

The effective exchange strength \(R_{\mathrm{joint}}=\sin^2\theta\) was
measured on a \(6\times6\) grid of intrinsic fractions \((f_A,f_B)\) and
compared with three composition-law hypotheses.

{\def\LTcaptype{none} % do not increment counter
\begin{longtable}[]{@{}lrl@{}}
\toprule\noalign{}
Hypothesis & Max error & Verdict \\
\midrule\noalign{}
\endhead
\bottomrule\noalign{}
\endlastfoot
Mean \((f_A+f_B)/2\) & \(7.2\times10^{-16}\) & \textbf{Exact
agreement} \\
Product \(f_Af_B\) & \(0.5\) & Rejected \\
\(\sqrt{f_Af_B}\) & \(0.5\) & Rejected \\
\end{longtable}
}

The coupling of the diagonal sector composes exactly additively (as the
mean). This is the grammar corresponding to the additivity of mass, and
no vertex product of \(q_Aq_B\) type exists in this sector (zeroth-order
no-go). No product structure was detected in the transfer observables
(swings of norm and localization) either (correlation \(|r|\le0.24\)).
Furthermore, the mean law does not depend on the readout convention:
under the variant masks (even≥6, any≥4) it holds identically for the
fractions measured per mask (max error \(4.4\times10^{-16}\); as a
theorem, it is a consequence of bin-wise power additivity).

\subsection{6. Experiment II: Charge Grammar of the Overlap Term
(Product, Sign,
Neutralization)}\label{experiment-ii-charge-grammar-of-the-overlap-term-product-sign-neutralization}

\subsubsection{6.1 Null Theorem (Carrier Separation = Closure of the
Coherent Flow
Path)}\label{null-theorem-carrier-separation-closure-of-the-coherent-flow-path}

In the standard configuration (with carrier), over all \(7\times7\)
pairs of A/B single harmonics,
\(|\langle A_j|B_k\rangle|\le2.7\times10^{-17}\) --- the two channels
separate into exactly orthogonal subspaces, and the overlap term is
identically zero. This is not a failure of measurement but a theorem: it
establishes that \textbf{the existence condition of charge-type
interaction is mode (channel) sharing}, and it provides the second
mechanism of neutralization (channel non-sharing).

\subsubsection{6.2 Record of the Invalid
Experiment}\label{record-of-the-invalid-experiment}

The carrierless + parity-reversed-mask configuration (v2) became
invalid: without the carrier, the harmonics do not land on the intended
bins and the fraction tuning collapses. This failure established that
the diagonal \(\theta\) readout (which presupposes the carrier) and the
coherent flow path (which presupposes sharing) sit on mutually exclusive
configurations in this toy model, and that a physical model requires
partial sharing (preserved as a record).

\subsubsection{6.3 Term-by-Term Verification of the Charge
Grammar}\label{term-by-term-verification-of-the-charge-grammar}

With carrierless anchors (shared channels measured:
\(c_0=c_6=0.500000\), cross-anchor leakage \(\le5\times10^{-16}\)) and
probe rotation angles \(\theta_p\in\{0.05,0.1,0.2,0.3\}\) (declared as
instrument settings), we measured the \(6\times6\) grid × 8 phase points
\(\varphi\) × equal/unequal norm (\(N_B=1.44\)). Against the exact
prediction of Section 3:

{\def\LTcaptype{none} % do not increment counter
\begin{longtable}[]{@{}
  >{\raggedright\arraybackslash}p{(\linewidth - 2\tabcolsep) * \real{0.5000}}
  >{\raggedright\arraybackslash}p{(\linewidth - 2\tabcolsep) * \real{0.5000}}@{}}
\toprule\noalign{}
\begin{minipage}[b]{\linewidth}\raggedright
Verification item
\end{minipage} & \begin{minipage}[b]{\linewidth}\raggedright
Result
\end{minipage} \\
\midrule\noalign{}
\endhead
\bottomrule\noalign{}
\endlastfoot
Modulation amplitude
\(M=\lvert\sin2\theta_p\rvert\sqrt{f_Af_B}\,\nu\,\lvert c_6\rvert\)
(vertex-product law) & Max error \(4.0\times10^{-15}\) (equal norm),
\(5.7\times10^{-15}\) (unequal norm) \\
Offset (magnitude term + bosonic cross term) & Max error
\(7.1\times10^{-15}\) (equal), \(6.5\times10^{-15}\) (unequal) \\
Neutrality (\(M=0\) at \(f_Af_B=0\)) & \(M\le4.2\times10^{-15}\) \\
Sign (flow direction reverses under a half turn of \(\varphi\)) &
Reversal in every cell where the modulation dominates the offset (63
equal / 76 unequal cells) \\
\end{longtable}
}

All three terms of the flow equation --- the diffusion term
(magnitude-driven), the bosonic cross term, and the charge term
(amplitude product × relative phase) --- matched the prediction at
machine precision. \textbf{The sign is the relative phase on the shared
channel, and the magnitude of the coupling is the product of
amplitudes.}

As a by-product, the shared channels have a ladder structure: \(A_k\)
couples only to \(B_{k\pm2}\), and all coupling values are exactly
\(0.5\) (selection rule \(\Delta k=\pm2\)).

\subsubsection{6.4 Claim-Hardening Experiments (Multiple Channels, Sign
Persistence, Partial Sharing, Distance
Law)}\label{claim-hardening-experiments-multiple-channels-sign-persistence-partial-sharing-distance-law}

Four additional experiments were carried out to close the unverified
aspects of the claims of the first draft (all runners and data are
committed).

\textbf{(i) Channel additivity of charge.} For mixed states carrying two
shared channels simultaneously (\(A_5\)--\(B_7\) and \(A_7\)--\(B_9\)),
the effective coupling agreed exactly with the \textbf{coherent sum of
per-channel amplitude products}

\[
\langle a|b\rangle=\sum_{\text{ch}}u\,v\,c_{\text{ch}}\,e^{i\varphi_{\text{ch}}}
\]

(\(6\times6\) two-phase grid, max error \(5.7\times10^{-15}\)). Charge
adds as amplitude across multiple channels.

\textbf{(ii) Multi-collision behavior of the sign.} Under iteration of a
fixed probe rotation, the transfer oscillates in Rabi fashion, and the
sign component (the component that reverses under phase reversal) is
exactly mirror-symmetric over the entire trajectory (antisymmetry error
\(2.2\times10^{-16}\)). The sign is therefore exact as the direction of
flow at each collision, but \textbf{speaking of a net steady flow
requires an argument about averaging over the oscillation} --- the same
situation in which the Bjerknes force is defined as a time average over
the acoustic period.

\textbf{(iii) Partially shared states = simultaneous validity of the two
grammars on the same state pair.} The exclusivity recorded in Section
6.2 (the diagonal readout requires the carrier; the coherent flow path
requires sharing) is resolved by the mixture

\[
a=\sqrt{1-w}\,A_{\mathrm{on}}+\sqrt{w}\,A_{\mathrm{off}},\qquad
b=\sqrt{1-w}\,B_{\mathrm{on}}+\sqrt{w}\,e^{i\varphi}B_{\mathrm{off}}
\]

of a carrier-on component and a carrierless shared component. On this
state pair, the diagonal readout \(R\) (bin-additive prediction, error
\(\le3.3\times10^{-16}\)) and the coherent modulation
\(M=\sin2\theta\cdot w\,|c|\) (error \(\le6.1\times10^{-16}\)) held
\textbf{simultaneously}. The coexistence of the two grammars is not a
switch between configurations but a property of a single state pair.

\textbf{(iv) Distance law of the coherent coupling.} The B wave packet
was translated exactly along \(\chi\) and the coupling
\(\langle a|b(d)\rangle\) was measured. The flow identity is exact at
all separations (\(\le1.4\times10^{-14}\)), but the distance law is
\textbf{entirely determined by the spectral envelope of the wave
packet}: for equally spaced comb states the coupling does not decay but
revives quasi-periodically (\(0.97\) even at \(d\ge32\)), while for
Gaussian-envelope packets it decays by a factor of about \(400\) over
the correlation length. In neither case does a \(1/d^2\) law appear.
\textbf{The original prediction of a ``short-range type'' was refuted}
(correction record (v)). As a consequence, the Coulomb long-range law
\(1/r^2\) cannot arise from direct overlap; a mediating mechanism
(harmonic closure {[}3{]}) must carry it --- this experimentally
pinpoints what the distance dictionary (Open Problem 2) must translate.

\subsubsection{6.5 Exclusion of the Reversed Assignment --- Uniqueness
of the Two-Grammar
Correspondence}\label{exclusion-of-the-reversed-assignment-uniqueness-of-the-two-grammar-correspondence}

With the measurements of Sections 3--6 in hand, the ``one-to-one
correspondence'' of the property table can be fixed as a corollary via
an exclusion argument. This section contains no new experiment --- what
follows is a logical consequence solely of rows already measured and
committed.

\textbf{Corollary (exclusion of the reversed assignment).} In this
system, the assignment of the magnitude term to the gravity-type grammar
and of the overlap term to the charge-type grammar is not an arbitrary
naming based on resemblance of properties but a unique correspondence
that excludes the reversed assignment.

\textbf{Proof.} The magnitude term \(\sin^2\theta\,(N_A-N_B)\) is a
function of norms alone and is identically invariant under independent
channel rephasings (Section 3, algebraic). Its composition law is
exactly additive (mean, Experiment I, error \(\le7.2\times10^{-16}\)),
and no product structure exists in the diagonal sector (zeroth-order
no-go). The magnitude term therefore possesses neither the two signs via
relative phase, nor flow reversal under phase reversal, nor an amplitude
product, nor neutralization via channel non-sharing, and \textbf{cannot
be the charge-type grammar}.

The overlap term \(\sin2\theta\,\mathrm{Re}\langle a|b\rangle\), on the
other hand, reads the amplitude product on shared modes (Section 6.3,
vertex-product law \(\le4.0\times10^{-15}\)), reverses sign under a half
turn of the relative phase (measured in Section 6.3; exactly
mirror-symmetric even over many collisions, Section 6.4 (ii),
\(2.2\times10^{-16}\)), and vanishes for zero components or non-shared
channels (neutrality \(\le4.2\times10^{-15}\); null theorem, Section
6.1). Reading it as the gravity-type grammar would make the sign of the
gravitational charge reverse under an internal phase flip and make
gravitational action vanish under orthogonalization / decoherence,
contradicting the requirements of the gravity-type grammar as a
phase-blind, universal readout of magnitude (universality, no sign,
unscreenable). \textbf{Hence the overlap term cannot be the gravity-type
grammar.}

Since, by the classification theorem of Section 4, no third term exists
in this class, no assignment other than

\[
\boxed{\text{magnitude term}=\text{gravity-type grammar},\qquad \text{overlap term}=\text{charge-type grammar}}
\]

exists. \(\blacksquare\)

What this corollary fixes is the \textbf{pre-geometric assignment of
interaction grammars}. The mapping to distance, direction, and metric on
the three emergent directions, and the identification with the physical
gravitational and Coulomb fields, remain as the separate tasks of the
scale dictionary (Open Problem 2) and the translation of sign (Open
Problem 3) --- this limitation of scope is not meant to weaken the claim
but to fix rigorously the domain of the uniqueness proposition.

\subsection{7. Experiment III: The Quantization Mechanism (Rational
Locking) --- Summary of Sister Paper
C}\label{experiment-iii-the-quantization-mechanism-rational-locking-summary-of-sister-paper-c}

This section is the summary, required for this paper's claims (the
quantization row of the property table), of the experiments reported in
full in sister paper C (lock dynamics).

\subsubsection{7.1 Zeroth-Order No-Go and the Working
Hypothesis}\label{zeroth-order-no-go-and-the-working-hypothesis}

By the theorem of Section 4, under a phase-blind readout \(\theta\) is a
constant of the motion (numerically confirmed: max drift
\(2.2\times10^{-16}\) over 400 collisions × 4 amplitudes), and neither
\(\theta\) dynamics nor locking exists. Introducing a phase-sensitive
readout (the power ratio of the symmetric channel \(X=A+B\)) as a
conservative working hypothesis makes \(\theta\) evolve.

\subsubsection{7.2 Measurement of the Arnold
Tongue}\label{measurement-of-the-arnold-tongue}

Sweeping the initial B amplitude, we found a finite interval where the
tail average of the late-time rotation number pins exactly to a
rational: for amplitudes \(2.52\)--\(3.3\),

\[
\overline{\theta/\pi}=1/3
\quad(\text{800 collisions, tail-300 average, }|{\overline{\theta/\pi}}-1/3|\le5.6\times10^{-17})
\]

\(\theta\) itself oscillates (tail standard deviation
\(\sim7\times10^{-2}\)), and only the mean rotation number locks to a
rational --- a periodic orbit exhibiting the standard behavior of an
Arnold tongue. In the root convention \(\theta=\pi/2-\pi m/n\) this
corresponds to \((m,n)=(1,6)\), an order-6 lock. Separately, with the
counting construction alone (zero amplitude search), an order-4 lock
(\(R=1/2\) exact, period-8 recurrence, integer reconstruction of
\((m,n)=(1,4)\)) was also obtained. The mechanism of quantization
plateaus really exists in an endogenous-coupling system whose coupling
is not supplied externally.

Note that the point at which the counting construction (equal-weight
points = rational \(R\)) can connect directly to rational-angle locking
(\(\theta/\pi\) rational) is, by Niven's theorem (when \(\theta/\pi\) is
rational, \(\cos\theta\) is rational only for
\(\cos\theta\in\{0,\pm\tfrac12,\pm1\}\) {[}17{]}), exactly limited to
the five points

\[
R \in \{0,\ 1/4,\ 1/2,\ 3/4,\ 1\}
\]

Of these, this paper demonstrated \(R=1/2\) (period 8). Verification of
the remaining nontrivial points \(R=1/4\) (predicted period 12) and
\(R=3/4\) (predicted period 6) is carried by sister paper B (under the
even≥2 variant of the convention audit, it has already been measured
that the counting state lands exactly on \(R=3/4\)).

\subsubsection{7.3 Record of the Limits of the
Search}\label{record-of-the-limits-of-the-search}

Just below the lower edge of the tongue, the rotation number varies
non-monotonically by \(\sim10^{-2}\) between adjacent amplitude points
(a symptom of the supercritical regime), so a naive sweep for high-order
tongues is unsuitable. The closest approach to the target
\(\rho=39/124\) was \(3.3\times10^{-4}\) (1500 collisions, scan width
\(0.002\)) with no plateau detected --- recorded as an upper bound on
the tongue width. In addition, the hypothesis of obtaining the winding
number \(m=23\) from a static state decomposition has been refuted
(winding-number spectroscopy experiment: the state spectrum agrees
completely with mapping artifacts; the \(k=23\) occupancy is at random
level).

\subsection{8. Convention Audit --- Summary of Sister Paper
B}\label{convention-audit-summary-of-sister-paper-b}

This section is the summary, required for the convention robustness of
this paper's claims, of the audit reported in full in sister paper B
(structure of counting readouts). The definition of the fermion mask was
varied over five variants (even≥4 / even≥6 / even≥2 / any≥4 /
parity-reversed odd≥3) to audit the convention dependence of the claims.

\textbf{Physics-grade (survives all conventions):} that the readout of
the equal-weight point is predicted by counting from the measured bin
table × mask (20/20, \(\le8\times10^{-16}\)); the small-denominator
rationality of the equal-weight points; the amplitude dependence; the
counting-to-lock beachhead (\(R=1/2\), period 8, \((1,4)\)). In
particular, under the parity-reversed mask the lock support moves to the
30 even harmonics --- \textbf{the assignment ``odd harmonics =
fermionic'' is a convention label, and the structure (counting → lock)
is independent of the label}.

\textbf{Bookkeeping-grade (products of the convention):} the concrete
form of the counting rule, the concrete values of the equal-weight
points, the assignment of the odd/even roles, the existence condition of
the beachhead. All claims of this paper are stated in physics-grade
form.

\subsection{9. Claims}\label{claims}

\textbf{Claim 1 (principal; discovery of the mechanism).} The two-body
flow of an anonymous closed wave system decomposes identically into two
terms, and no third term exists. The two terms read exclusively the two
pieces of information in the wave shape (norm / relative phase on shared
modes), and interaction necessarily differentiates into exactly two
grammars:

{\def\LTcaptype{none} % do not increment counter
\begin{longtable}[]{@{}
  >{\raggedright\arraybackslash}p{(\linewidth - 4\tabcolsep) * \real{0.3333}}
  >{\raggedright\arraybackslash}p{(\linewidth - 4\tabcolsep) * \real{0.3333}}
  >{\raggedright\arraybackslash}p{(\linewidth - 4\tabcolsep) * \real{0.3333}}@{}}
\toprule\noalign{}
\begin{minipage}[b]{\linewidth}\raggedright
\end{minipage} & \begin{minipage}[b]{\linewidth}\raggedright
Magnitude term
\end{minipage} & \begin{minipage}[b]{\linewidth}\raggedright
Overlap term
\end{minipage} \\
\midrule\noalign{}
\endhead
\bottomrule\noalign{}
\endlastfoot
Composition law & \textbf{Additive} (mean, \(\le7\times10^{-16}\)) &
\textbf{Product} \(\sqrt{f_Af_B}\) (\(\le4\times10^{-15}\)) \\
Sign & None & Present (reverses with relative phase, measured) \\
Screening & Impossible (always on) & Possible (two mechanisms: zero
component / non-sharing) \\
Quantization & None (continuous) & Present (rational locking,
measured) \\
Range of the coupling & Ratio \((R/R_0)^2\): unbounded, continuous &
Angle: \(\alpha\le1/4\pi\) and lock discretization confine it to the
\(O(10^{-1})\) band \\
\end{longtable}
}

This property table corresponds one-to-one with the contrast table of
gravity and electromagnetism. Moreover, this correspondence is
\textbf{uniquely determined by an exclusion argument}, since both
directions of the reversed assignment (overlap term → gravity-type /
magnitude term → charge-type) contradict properties already measured
(Section 6.5, Corollary). The differentiation of the forces requires no
external input of particle species, coupling constants, or background.

\textbf{Claim 2 (completeness on the Coulomb side).} The overlap term
possesses all five elements of the Coulomb structure --- inverse square
(derived in {[}3{]}), product, sign, quantization, and value (the exact
finite-order root \(\sin^2(23\pi/124)\), matching \(\sqrt{4\pi\alpha}\)
to 0.16 ppm {[}2{]}). We call this quantity \textbf{the model's
elementary charge}.

\textbf{Claim 3 (hierarchy).} A long-range interaction that is weaker by
orders of magnitude can exist only in the grammar of the magnitude term
(the overlap term, squeezed between the unitarity bound and
quantization, cannot leave the angular band). Moreover, the coupling of
the magnitude term is of scale-ratio type \((R/R_0)^2\), and in any
universe containing localized closures \(R\ll R_0\); that is,
\textbf{the enormous hierarchy between the two great long-range forces
is a structure that follows from the definition of localization ---
particles are far smaller than the background}. The premise
(gravitational coupling = norm ratio) and the derivation of the
hierarchy's exponent (the scale dictionary) are left as explicit tasks.

\textbf{Claim 4 (the one remaining item).} Toward identification with
the physical elementary charge, the only work remaining inside the model
is the derivation of the mechanism that selects order 124 (the overall
picture of the tasks connecting to emergent spacetime and physical
forces is organized, as a series separate from the internal selection
problem, in the closing section of sister paper C). The candidates have
been narrowed by elimination through the experiments of sister paper C:
the static decomposition does not select (refuted); the fixed-point-type
channel shows no rational concentration (eliminated); the tongue weights
of locking collapse with the denominator at a measured ratio \(>461\)
(insufficient); and observational selection shows fidelity 1 at exact
roots together with \textbf{inheritance of the selection into the
divisor structure of the observation clock} (an observer with clock
\(J=8\) cannot see the order-124 recurrence, while one with \(J=248\)
can). The selection problem has therefore been transformed into the
single falsifiable question \textbf{``why is the observation clock
commensurable with 248?''}.

\textbf{The sign-discrimination problem (prediction).} The sign was
measured as a relative phase. Its translation into spatial
attraction/repulsion requires the distance dictionary
(\(\Delta\theta\leftrightarrow r\), unfinished), but the moment the
dictionary is completed, whether this system is Bjerknes-type (equal
phase = attraction) or electromagnetic-type (equal sign = repulsion) is
uniquely determined by the structure of the dictionary --- the sign
problem is not an ambiguity but an automatically verified prediction.

\subsection{10. Convergence with Known
Systems}\label{convergence-with-known-systems}

The force law product × cos(relative phase) × inverse square is the
structure that Bjerknes (1875) demonstrated both theoretically and
experimentally with pulsating spheres in an ideal fluid (alive in modern
acoustics as the secondary Bjerknes force); that mode overlap is the
existence condition of coupling is realized by visibility and
distinguishability in optics (including Hong--Ou--Mandel); quantization
plateaus via phase locking by the Shapiro steps of superconductivity
(and, recently, current quantization via Bloch-oscillation
synchronization); and the correspondence between the force ratio
\(\sim10^{40}\) and scale ratios by Dirac's large numbers (1937) ---
each independently realized and observed. That fluids, optics, and
superconductivity --- mutually unrelated --- and this model converge on
the same grammar is evidence that this grammar is a general structure of
closed wave systems rather than a detail of the medium.

The term ``emergent electromagnetism'' also appears in recent
canonical-theoretic work (Phys. Rev.~D \textbf{110}, 125006, 2024), but
the distinction made there (the canonical structure of
fundamental/emergent fields) and the distinction made here (the
exclusive readout of the two terms of the flow) are different things.
The contributions specific to this paper are: (i) the classification
theorem for readouts and the composition laws in an endogenous-coupling
system in which the coupling is read from the state rather than being an
instrument constant; (ii) that the five elements and the two grammars
hold simultaneously as terms of a single identity of a single system;
(iii) the machine-precision verification of every row together with the
convention audit.

\subsection{11. Open Problems}\label{open-problems}

\begin{enumerate}
\def\labelenumi{\arabic{enumi}.}
\tightlist
\item
  \textbf{Selection problem}: identification of the endogenous dynamics
  that selects order 124. The candidates have been greatly narrowed by
  the experiments of sister paper C --- the fixed-point-type channel is
  eliminated (the attractor set is a continuous band, with no rational
  concentration); by the measured tongue-width ratio (denominator 3 vs
  denominator ≥4, \(>461\)) the lock weights are insufficient;
  observational selection shows fidelity 1 at exact roots and
  \textbf{inheritance of the selection into the divisor structure of the
  observation clock}, but finite-count observation can narrow only down
  to the recurrence basin. The current form of the question is therefore
  \textbf{``why is the observation clock commensurable with 248?''}, and
  this final form is taken up by the parent paper (the main paper 9)
\item
  \textbf{Scale dictionary}: the distance translation
  \(\Delta\theta_n\leftrightarrow r\) and the value of \(R/R_0\) (the
  exponent of the hierarchy). Section 6.4 (iv) identified what the
  dictionary must translate: since direct overlap does not produce
  \(1/r^2\), the dictionary must translate \textbf{the distance readout
  of the mediating mechanism (harmonic closure)}
\item
  \textbf{Translation of sign}: discrimination between Bjerknes-type and
  electromagnetic-type (decided automatically once the dictionary is
  complete)
\item
  \st{Partial-sharing model} \textbf{Solved} (Section 6.4 (iii):
  demonstrated that the two grammars hold simultaneously on the same
  state pair)
\item
  \textbf{General form of the counting rule}: the closed form of the
  counting structure of equal-weight points and the general theory of
  integer-pair readout \textbf{are carried by sister paper B}
\end{enumerate}

\subsection{12. Reproducibility}\label{reproducibility}

All experiments referenced a single core runner (unchanged \(\theta\)
readout and rotation) with SHA-256 verification, were conducted by
writing predicted values as in-code constants before measurement, and by
placing reproductions of known results (anchors) at the head of each
runner. Refuting experiments, invalid experiments, and corrections
(including the correction of the R4 conservation misidentification) were
all recorded without deletion. Experiment folders and commits:

{\def\LTcaptype{none} % do not increment counter
\begin{longtable}[]{@{}
  >{\raggedright\arraybackslash}p{(\linewidth - 4\tabcolsep) * \real{0.3333}}
  >{\raggedright\arraybackslash}p{(\linewidth - 4\tabcolsep) * \real{0.3333}}
  >{\raggedright\arraybackslash}p{(\linewidth - 4\tabcolsep) * \real{0.3333}}@{}}
\toprule\noalign{}
\begin{minipage}[b]{\linewidth}\raggedright
Experiment
\end{minipage} & \begin{minipage}[b]{\linewidth}\raggedright
Folder
\end{minipage} & \begin{minipage}[b]{\linewidth}\raggedright
Commit
\end{minipage} \\
\midrule\noalign{}
\endhead
\bottomrule\noalign{}
\endlastfoot
Two-memory discrimination battery &
\texttt{two\_memory\_readout\_battery\_pre\_v1} & 5444fbdb \\
Counting-rule verification + direct counting-lock bridge &
\texttt{counting\_rule\_lock\_bridge\_pre\_v1} & 8dff131e \\
Winding-number spectroscopy (refutation) &
\texttt{winding\_spectroscopy\_pre\_v1} & 641048ad \\
Convention audit & \texttt{convention\_audit\_pre\_v1} & 5a8f4b0a \\
Mode-locking probe (Z₆ discovery) &
\texttt{mode\_locking\_probe\_pre\_v1} & 56453653 \\
Readout classification + tongue spectroscopy &
\texttt{tongue\_spectroscopy\_pre\_v1} & 913756fe \\
Mixed coupling (mean law) & \texttt{mixed\_coupling\_pre\_v1} &
fd7a3745 \\
Cross-term charge grammar (v1 null / v2 invalid / v3 established) &
\texttt{cross\_term\_charge\_pre\_v1} & 26f26eb8 \\
Batch regeneration of all figures in this paper &
\texttt{paperA\_two\_grammar\_evidence\_v1} & 381c2c5c \\
Niven intersection locks (periods 6/8/12, all predictions confirmed) &
\texttt{niven\_points\_lock\_pre\_v1} & 8f99d5c6 \\
Claim hardening (multiple channels / sign persistence / mean-law
convention extension) & \texttt{claim\_hardening\_battery\_v1} &
47e2da02 \\
Partial sharing + fixed-point map &
\texttt{partial\_sharing\_fixed\_point\_v1} & 47e2da02 \\
Distance law (comb/Gaussian contrast) & \texttt{coupling\_range\_v1} &
47e2da02 \\
Observation selection (clock selectivity) &
\texttt{observation\_selection\_v1} & 47e2da02 \\
Tongue-width measurement & \texttt{tongue\_width\_measurement\_v1} &
47e2da02 \\
\end{longtable}
}

Figures (all regenerable by the single runner of
\texttt{paperA\_two\_grammar\_evidence\_v1}):

\begin{itemize}
\tightlist
\item
  Fig. A \texttt{figA\_carrier\_overlap\_v1}: overlap matrices (with
  carrier = all zeros / without = ladder \(\Delta k=\pm2\))
\item
  Fig. B \texttt{figB\_classification\_v1}: conserved/dynamical
  classification of the 7 readouts
\item
  Fig. C \texttt{figC\_mean\_law\_v1}: mean composition law and
  residuals
\item
  Fig. D \texttt{figD\_charge\_grammar\_v1}: sign reversal of
  \(dN_B(\varphi)\) and the vertex-product line
\item
  Fig. E \texttt{figE\_three\_terms\_v1}: three-term decomposition,
  prediction vs.~measurement (equal/unequal norm)
\end{itemize}

All measured-row CSVs were bundled by explicitly overriding the
repository's exclusion rules.

\begin{center}\rule{0.5\linewidth}{0.5pt}\end{center}

\section{References}\label{references}

\subsection{Self-citations}\label{self-citations}

\begin{enumerate}
\def\labelenumi{\arabic{enumi}.}
\tightlist
\item
  Noriaki Kihara, ``Anonymous Equal-Amplitude Composite-Wave Model:
  Basic Axiom System v6'', Version DOI:
  \texttt{10.5281/zenodo.21465984}, Concept DOI:
  \texttt{10.5281/zenodo.21315735}, 2026.
\item
  Noriaki Kihara, ``Discovery of Finite-Order Resonance in Iterated
  Exchange Scattering --- Identifying the Origin of Peaks near
  Fine-Structure-Constant Values 137 and 128 with a Reproducible
  Wave-Packet Model'', Version DOI: \texttt{10.5281/zenodo.21421367},
  Concept DOI: \texttt{10.5281/zenodo.21421366}, 2026. (The two-channel
  exchange operator \(U_R\) and the eigenvalue structure
  \(\lambda_a=e^{-2\pi i m/n}\), the B63 construction, and the defining
  source of the exact finite-order roots)
\item
  Noriaki Kihara, ``Future Phase-Position Acceleration Map and the
  Inverse-Square Law via Harmonic Closure in an AB Two-Body Closed Phase
  System v4'', Version DOI: \texttt{10.5281/zenodo.21468270}, Concept
  DOI: \texttt{10.5281/zenodo.21441081}, 2026. (Derivation of the
  inverse-square law, and the precedent for the three-step procedure
  zeroth-order no-go → conservative working hypothesis → independent
  verification)
\item
  Noriaki Kihara, ``Where Are Mass, Momentum, and Energy Read From? ---
  Multigauge Conserved Readouts of the ABC Closed Phase System'',
  Version DOI: \texttt{10.5281/zenodo.21308050}, Concept DOI:
  \texttt{10.5281/zenodo.21308049}, 2026.
\item
  Noriaki Kihara, ``Preliminary Summary of Acceleration-Basis and
  Localization Exchange in Exchange-Interference Scattering-Matrix
  Fermion-Like Collisions v1'', Version DOI:
  \texttt{10.5281/zenodo.21333768}, Concept DOI:
  \texttt{10.5281/zenodo.21333766}, 2026. (Code origin of the scattering
  engine)
\item
  Noriaki Kihara, ``Selection of the State-Exchange Weight G\_R = 1 − R
  and Candidate Correspondence with the Fine-Structure Constant: A
  Numerical Experiment v1'', Version DOI:
  \texttt{10.5281/zenodo.21396761}, Concept DOI:
  \texttt{10.5281/zenodo.21396760}, 2026. (Origin of the System A code)
\end{enumerate}

\begin{itemize}
\tightlist
\item
  Sister paper B: Noriaki Kihara, ``Structure of Counting Readouts in an
  Anonymous Two-Channel Closed Wave System'' (Part II of the trilogy),
  Version DOI: \texttt{10.5281/zenodo.21763998}, Concept DOI:
  \texttt{10.5281/zenodo.21763997}, 2026.
\item
  Sister paper C: Noriaki Kihara, ``Lock Dynamics in an Anonymous
  Two-Channel Closed Wave System'' (Part III of the trilogy), Version
  DOI: \texttt{10.5281/zenodo.21764000}, Concept DOI:
  \texttt{10.5281/zenodo.21763999}, 2026.
\end{itemize}

\subsection{External References}\label{external-references}

The external references are not the basis of the derivations in this
paper. They are used to show independent realizations of the same
structure (evidence of convergence) and historical precedence.

\begin{enumerate}
\def\labelenumi{\arabic{enumi}.}
\setcounter{enumi}{6}
\tightlist
\item
  G. Forbes, ``Hydrodynamic Analogies to Electricity and Magnetism'',
  \emph{Nature} \textbf{24}, 360--361 (1881). DOI:
  \texttt{10.1038/024360a0}. (Report of C. A. Bjerknes's
  pulsating-sphere apparatus)
\item
  L. A. Crum, ``Bjerknes forces on bubbles in a stationary sound
  field'', \emph{J. Acoust. Soc. Am.} \textbf{57}, 1363--1370 (1975).
  DOI: \texttt{10.1121/1.380614}. (The secondary Bjerknes force:
  amplitude product × cos(phase difference) × inverse square; in-phase
  attraction, antiphase repulsion)
\item
  M. Born and E. Wolf, \emph{Principles of Optics}, 7th ed., Cambridge
  University Press (1999). (Two-beam interference law, mutual coherence)
\item
  C. K. Hong, Z. Y. Ou, and L. Mandel, ``Measurement of subpicosecond
  time intervals between two photons by interference'', \emph{Phys.
  Rev.~Lett.} \textbf{59}, 2044 (1987). DOI:
  \texttt{10.1103/PhysRevLett.59.2044}.
\item
  L. Mandel, ``Coherence and indistinguishability'', \emph{Opt. Lett.}
  \textbf{16}, 1882 (1991). DOI: \texttt{10.1364/OL.16.001882}. (Mode
  distinguishability destroys interference visibility = coherent
  coupling)
\item
  S. Shapiro, ``Josephson currents in superconducting tunneling: the
  effect of microwaves and other observations'', \emph{Phys. Rev.~Lett.}
  \textbf{11}, 80 (1963). DOI: \texttt{10.1103/PhysRevLett.11.80}.
  (Quantization plateaus via phase locking)
\item
  M. H. Jensen, P. Bak, and T. Bohr, ``Complete devil's staircase,
  fractal dimension, and universality of mode-locking structure in the
  circle map'', \emph{Phys. Rev.~Lett.} \textbf{50}, 1637 (1983). DOI:
  \texttt{10.1103/PhysRevLett.50.1637}. (Arnold tongues, devil's
  staircase)
\item
  P. A. M. Dirac, ``The cosmological constants'', \emph{Nature}
  \textbf{139}, 323 (1937). DOI: \texttt{10.1038/139323a0}.
  (Large-number correspondence between force ratios and scale ratios.
  The same hypothesis's prediction of a time-varying G is not forced by
  the framework of this paper)
\item
  O. Klein, ``Quantentheorie und fünfdimensionale Relativitätstheorie'',
  \emph{Z. Phys.} \textbf{37}, 895--906 (1926). DOI:
  \texttt{10.1007/BF01397481}.
\item
  E. I. Duque, ``Emergent electromagnetism'', \emph{Phys. Rev.~D}
  \textbf{110}, 125006 (2024). DOI:
  \texttt{10.1103/PhysRevD.110.125006}. arXiv:\texttt{2407.14954}. (The
  term is close, but it concerns the fundamental-field/emergent-field
  distinction via canonical structure, distinct from the two-term
  decomposition of this paper)
\item
  I. Niven, \emph{Irrational Numbers}, Carus Mathematical Monographs
  No.~11, Mathematical Association of America (1956). (Theorem on the
  rationality of cosines of rational angles; the basis for the
  intersection set \(\{0,1/4,1/2,3/4,1\}\) of the counting
  rational-\(R\) family and the rational-angle lock family)
\end{enumerate}

\begin{center}\rule{0.5\linewidth}{0.5pt}\end{center}

\textbf{Addendum (record of corrections)} The following errors occurred
and were corrected in the course of this work. (i) The cross-spectrum
readout R4 was misidentified as a conserved readout (tail freezing
misread as conservation) → the classification experiment revealed
fixed-point convergence, and it was corrected. (ii) A mix-up in the sign
prediction of the offset term of the flow equation → resolved by
calibrating the rotation convention (\(\sigma=+1\)) (the exact agreement
of the error \(0.389=\sin0.4\) was the evidence of the sign slip). (iii)
Collapse of the fraction-tuning assumption in the carrierless
configuration (v2 invalid). (iv) The error in a prior record stating
that the intersection of the counting-rational family and the
rational-angle lock family is ``\(R=1/2\) only'' → corrected by Niven's
theorem to the five points \(\{0,1/4,1/2,3/4,1\}\) (the correction was
appended to the experiment folder README, and the remaining two points
were confirmed with all predictions confirmed in the Niven intersection
experiment). (v) The erroneous prediction that the distance law of the
coherent coupling is ``short-range type'' → measurement revealed
undamped revival for comb states and envelope decay for Gaussian
packets; the distance law is determined by the packet envelope (Section
6.4 (iv)). All are recorded in the main text and in the reproduction
package.

\end{document}
